%!TEX root = ../../../main.tex

\chapter{RAG- und Retrieval-Konzept}
\label{chap:retrieval}

\paragraph{Einbettung.} Vor jedem Chunk steht der Hierarchie-Präfix \texttt{[Ordner | Datei | Kapitel]}. So wird \enquote{34,50 ct/kWh} als Arbeitspreis \emph{im Stromvertrag} eingebettet. Das Modell ist nomic-embed-text (768 Dimensionen, 8k Kontext \cite{nussbaum2024}), mit den Präfixen \texttt{search\_document:} und \texttt{search\_query:}. Dokumente und Anfragen laufen durch dieselbe Funktion. Modell und Präfixe stehen als Signatur neben dem Index und werden bei jeder Abfrage geprüft; Version1 lieferte bei einem anderen Embedding-Modell stillschweigend sinnlose Treffer.

\paragraph{Suche und Ablation.} ChromaDB liefert 20 Kandidaten, danach gilt $\text{final} = \text{base} \times \text{boost}$ mit $\text{base} = 1 - d_{\cos}$. Der Boost greift, wenn ein Wort der Frage zu Ordner, Datei oder \texttt{heading\_stack} passt. Wegen deutscher Komposita vergleiche ich über Präfixe, aber \enquote{Heizstrom} passt bewusst nicht zu \enquote{Strom}. \gls{BM25} \cite{robertson2009}, ein Cross-Encoder-Reranker \cite{nogueira2019} und Query-Rewriting sind Schalter. Sie gehen nur an, wenn sie messbar besser sind: Die mittlere \gls{MRR}-Differenz muss mindestens $+0{,}01$ betragen und das 95-\%-Bootstrap-Intervall über null liegen. Auf 93 Entwicklungsfragen (Tabelle~\ref{tab:ablation} im Anhang) bringt der Boost 1,5 nur $+0{,}011$ gegenüber keinem Boost, das Intervall berührt null. 1,25 und 2,0 ergeben dieselbe MRR, also bleibt die Vorgabe. BM25 ($+0{,}010$) und Rewriting per Regeln ($+0{,}011$) liegen im Messfehler. BM25 senkt zudem Beleg@5 von 88 auf 66\,\%. Der Reranker verliert 0,045 und kostet 1,1\,s, LLM-Rewriting verliert 0,37 und kostet 2,5\,s. Recall@5 liegt schon ohne Zusätze bei 100\,\%: Die Suche ist nicht der Engpass.

\paragraph{GraphRAG (Leitfrage~5).} \emph{Entitäten} wären Kunde, Vertragspartner, Vertrag, Rechnung, Zähler und Tarif. \emph{Beziehungen} wären \enquote{Rechnung gehört zu Vertrag}, \enquote{Kunde hat Vertrag} und \enquote{Vertrag läuft über Zähler}. \emph{Besser beantwortbar} wären Multi-Hop- und Überblicksfragen \cite{edge2024}. \emph{Klassisches RAG mit Kerndaten reicht}, weil die einzige tragende Beziehung ein Gleichheitsvergleich auf Kunden- oder Vertragsnummer ist. Bei 16 Dokumenten ist das ein Join über eine Spalte, den \texttt{get\_key\_facts} und \texttt{list\_contracts} liefern. \emph{Die Kosten} eines Graphen wären eine LLM-Extraktion je Chunk. Version1 hatte dafür 1326 Zeilen Code und fünf Werkzeuge, ohne den Nutzen je zu messen. Ich habe den Join gemessen: Auf acht Multi-Hop- und Listenfragen löst P2 mit beiden Werkzeugen 4, ohne sie 2. Gewonnen hat P2 die Vergleiche über Dokumentgrenzen. Die übrigen Fehler sind Lesefehler, keine fehlende Graphstruktur. Bei acht Fragen ist das Intervall weit (50\,\% [22--78] gegen 25\,\% [7--59]). Die Richtung stützt dennoch die Entscheidung gegen einen Graphen.
